\documentclass[a4paper, 12pt]{extarticle}


%%%%%%%%%%%%%%%%%%%%%%%%%%%%%%
% PACKAGES
%%%%%%%%%%%%%%%%%%%%%%%%%%%%%%

%!TEX encoding = UTF8

%%%%%%%%%%%%%%%%%%%%%%%%%%%%%%%%%
% PACKAGE IMPORTS
%%%%%%%%%%%%%%%%%%%%%%%%%%%%%%%%%


\usepackage[french]{babel}
\usepackage[T1]{fontenc}
\usepackage[utf8]{inputenc}

\usepackage[
	margin=2cm,
	right=2.2cm,
	footskip=1.5cm,
]{geometry}

%ams
\usepackage{amsmath,amsfonts,amsthm,amssymb,mathtools}

\usepackage{bookmark}
\usepackage{enumitem}
\usepackage[most,many,breakable]{tcolorbox}
\usepackage{varwidth,etoolbox}
\usepackage[makeroom]{cancel}
\usepackage{xcolor}
\usepackage{multicol,array}
\usepackage[ruled,vlined,linesnumbered]{algorithm2e}

\usepackage{caption, subcaption}

\usepackage{marginnote}

\usepackage{makecell} % pour \thead 

%index
\usepackage[texindy]{imakeidx}
\makeindex[
	columns=2, 
	title=Index, 
	options= -L french,
	%intoc,
]
\indexsetup{level=\chapter}

%virgules
\usepackage{icomma}

% for striked out implies sign (\centernot\implies)
\usepackage{centernot}

% roman numerals for \section
\renewcommand{\thesection}{\Roman{section}} 

% tikz
\usepackage{tikz}

% package systeme pour les systèmes d'équations bien alignés
\usepackage{systeme}
%\sysalign{r,r}
%\syseqspace{3pt}
%\syssignspace{3pt}

\newcommand{\sys}[2]{\systeme{#1{,}, #2.}}

% date
\usepackage{advdate}

\pagestyle{fancy}
\fancyhead[L]{Université de Mayotte}
\fancyhead[C]{\bf Série 1}
\fancyhead[R]{\today}

%%%%%%%%%%%%%%%%%%%%%%%%%%%%%%
% SELF MADE COLORS
%%%%%%%%%%%%%%%%%%%%%%%%%%%%%%

\input{../../../colors}

%%%%%%%%%%%%%%%%%%%%%%%%%%%%
% LETTERFONTS
%%%%%%%%%%%%%%%%%%%%%%%%%%%%

\input{../../../letterfonts}

%%%%%%%%%%%%%%%%%%%%%%%%%%%%
% MACROS
%%%%%%%%%%%%%%%%%%%%%%%%%%%%

%!TEX encoding = UTF8


%%%%%%%%%%%%%%%%%%%%%%%%%%%%
% THEOREMS
%%%%%%%%%%%%%%%%%%%%%%%%%%%%

\newcommand{\thm}[3]{\begin{theorem}[#1]\label{#3}#2\end{theorem}}
\newcommand{\cor}[3]{\begin{corollaire}[#1]\label{#3}#2\end{corollaire}}
\newcommand{\lem}[3]{\begin{lemme}[#1]\label{#3}#2\end{lemme}}
\newcommand{\prop}[3]{\begin{proposition}[#1]\label{#3}#2\end{proposition}}
\newcommand{\ex}[3]{\begin{exemple}[#1]\label{#3}#2\end{exemple}}
\newcommand{\dfn}[3]{\begin{definition}[#1]\label{#3}#2\end{definition}}
\newcommand{\qs}[2]{\begin{question}[#1]#2\end{question}}

\renewcommand\qedsymbol{$\blacksquare$}
\newcommand{\pf}[2]{\begin{preuve}[#1]#2\hfill\qedsymbol\end{preuve}}

\newcommand{\nt}[1]{\begin{remarque}#1\end{remarque}}
\newcommand{\str}[1]{\begin{strategie}#1\end{strategie}}
\newcommand{\mth}[1]{\begin{methode}#1\end{methode}}
\newcommand{\ax}[3]{\begin{axiome}[#1]\label{#3}#2\end{axiome}}
\newcommand{\notations}[1]{\begin{notation}#1 \end{notation}}
\newcommand{\nomen}[1]{\begin{nomenclature}#1 \end{nomenclature}}
\newcommand{\heur}[1]{\begin{heuristique}#1\end{heuristique}}
\newcommand{\motiv}[1]{\begin{motivation}#1\end{motivation}}


%%%%%%%%%%%%%%%%%%%%%%%%%%%%
% EXERCISES
%%%%%%%%%%%%%%%%%%%%%%%%%%%%

\newcommand{\exe}[4]{
	\begin{Exercise}[#1, label=#3]
		%\marginpar{\mbox{\scriptsize(solution p.\pageref{\ExerciseLabel-Answer})}}
		#2
	\end{Exercise}
	\begin{Answer}[ref=#3]
		#4
	\end{Answer}
}

% version 1
\newcommand{\exemulticols}[5]{
	\begin{multicols}{2}
	\begin{Exercise}[#1, label=#4]
		#2
	\end{Exercise}
	\columnbreak
		#3
	\end{multicols}
	\begin{Answer}[ref=#4]
		#5
	\end{Answer}
}

% version 2
%\newcommand{\exemulticols}[5]{
%	\begin{Exercise}[#1, label=#4]
%	\begin{multicols}{2}
%		#2
%	\columnbreak
%		#3
%	\end{multicols}
%	\end{Exercise}
%	\begin{Answer}[ref=#4]
%		#5
%	\end{Answer}
%}

\newcommand{\exeTF}[1]{}

%%%%%%%%%%%%%%%%%%%%%%%%%%%%
% OPERATORS
%%%%%%%%%%%%%%%%%%%%%%%%%%%%

% deliminators
\DeclarePairedDelimiter{\abs}{\lvert}{\rvert}
\DeclarePairedDelimiter{\norm}{\lVert}{\rVert}
\DeclarePairedDelimiter{\scal}{\langle}{\rangle}

\DeclarePairedDelimiter{\ceil}{\lceil}{\rceil}
\DeclarePairedDelimiter{\floor}{\lfloor}{\rfloor}
\DeclarePairedDelimiter{\round}{\lfloor}{\rceil}

%  - others
\DeclareMathOperator{\Lap}{\mathcal{L}}
\DeclareMathOperator{\Var}{Var} % variance
\DeclareMathOperator{\Cov}{Cov} % covariance
\DeclareMathOperator{\proj}{proj} % projection
\DeclareMathOperator{\sym}{sym} % symétrique

\DeclareMathOperator{\RE}{Re} % partie réelle
\DeclareMathOperator{\IM}{Im} % partie imaginaire

\DeclareMathOperator{\pgcd}{pgcd}
\DeclareMathOperator{\ppcm}{ppcm}

\DeclareMathOperator{\ev}{ev}

\DeclareMathOperator{\tr}{tr}
\DeclareMathOperator{\sgn}{sgn}

\DeclareMathOperator{\Frac}{Frac}

% Lois
\DeclareMathOperator{\Bern}{\text{Bern}}
\DeclareMathOperator{\Binom}{\text{Binom}}

% I prefer the slanted \leq
\let\oldleq\leq % save them in case they're every wanted
\let\oldgeq\geq
\renewcommand{\leq}{\leqslant}
\renewcommand{\geq}{\geqslant}


%%%%%%%%%%%%%%%%%%%%%%%%%%%%
% commands
%%%%%%%%%%%%%%%%%%%%%%%%%%%%

% tel que
\newcommand{\tqs}{\text{ tels que }}
\newcommand{\tq}{\text{ tq. }}
\newcommand{\et}{\text{ et }}
\newcommand{\ou}{\text{ ou }}
\newcommand{\pourtout}{\text{ pour tout }}
\newcommand{\sct}{\text{ sachant }}
\newcommand{\id}{\text{id}}
\newcommand{\apriori}{\emph{a priori}}


% ensemble avec bigl et bigr
\newcommand{\bigset}[1]{\bigl\{ #1 \bigr\}}
\newcommand{\Bigset}[1]{\Bigl\{ #1 \Bigr\}}
\newcommand{\bigpar}[1]{\bigl( #1 \bigr)}
\newcommand{\Bigpar}[1]{\Bigl( #1 \Bigr)}

% PLUS INFTY AND MINUS INFTY WITH NO SPACE
\newcommand{\pinfty}{{+}\infty}
\newcommand{\minfty}{{-}\infty}

% vecteur flèche
\renewcommand{\vec}[1]{\overrightarrow{#1}}

% vecteur pmatrix
\newcommand{\pvec}[2]{\begin{pmatrix} #1 \\ #2 \end{pmatrix}}

% vecteur norme
%\newcommand{\norm}[1]{\left\Vert #1 \right\Vert}

% point plan
\newcommand{\point}[3]{
	#1\left(#2 ; #3 \right)
}

% \smash avant \underline pour coller la ligne au mot
% attention : ça empêche les line breaks donc c'est pas super...
\let\oldunderline\underline
\renewcommand{\underline}[1]{\oldunderline{\smash{#1}}}

% emph + index
% i add underline just for now
\newcommand{\emphindex}[1]{\emph{\uline{#1}}\index{#1}}

% tableau croisé
\newcommand{\tableaucroise}[4]{
\begin{tabular}{|c|c|c|c|}
	\cline{2-4}
	\multicolumn{1}{c|}{} & #1 \\ \hline
	#2 \\ \hline
	#3  \\ \hline
	#4  \\ \hline
\end{tabular}
}

% emdash in mathmode
\newcommand{\dash}{\text{---}}

%%%%%%%%%%%%%%%%%%%%%%%%%%%%
% ENVIRONMENTS
%%%%%%%%%%%%%%%%%%%%%%%%%%%%

% two columns toc
\newcommand{\twocolstableofcontents}{
	\renewcommand*\contentsname{}
	\begin{center}  
		{\Large \bfseries TABLE DES MATIÈRES \bigskip}    
	\end{center}
	\hrule
	\begin{multicols}{2}\setlength{\columnseprule}{1pt}
		\titleformat{\chapter}{}{}{0pt}{} 
		\titlespacing*{\chapter}{0pt}
		{*7} %  % vertical space before the title <<<<<<<<<<
		{*-15}  %  idem after title (in ex units + glue) <<<<<<<<<<<
		\tableofcontents
	\end{multicols}
	\hrule
}

% python minted
\newcommand{\python}[1]{
\inputminted[
		linenos,
		gobble=0,
		breaklines=true, % otherwise it breaks for no apparent reason?
		breakafter=,,
		fontsize=\normalsize,
		numbersep=8pt,
		tabsize=4, % tab ident = 4 spaces
		fontfamily=courier, %important pour les signes <, >
]{python}{python/#1.py}
}



%%%%%%%%%%%%%%%%%%%%%%%%%%%%
% PACKAGE-SPECIFIC ENVIRONMENTS
%%%%%%%%%%%%%%%%%%%%%%%%%%%%

% global settings for the first level in enumerate environments
\setlist[enumerate, 1]{%
	align = left, labelwidth = \parindent%, labelsep = 0pt, leftmargin = \parindent
}


%%%%%%%%%%%%%%%%%%%%%%%%%%%%%%
% EXERCISES
%%%%%%%%%%%%%%%%%%%%%%%%%%%%%%

\usepackage{ifthen}
\renewcommand{\ExerciseHeader}{
	\tikz[baseline=(R.base)]\node[draw,rectangle, thick, inner sep=2pt](R) {\textbf{\hyperref[\ExerciseLabel-Answer]{\theExercise}.}};\!
	\ifnum\ExerciseDifficulty=0
	\else
		(\theExerciseDifficulty)
	\fi
}
\renewcommand{\DifficultyMarker}{$\star$}
\renewcommand{\AnswerHeader}{
	\tikz[baseline=(R.base)]\node[draw,rectangle, thick, inner sep=2pt](R) {\textbf{\theExercise.}};\!
}
\renewcommand{\exe}[4]{
	\begin{Exercise}[#1, label=#3]
		\if\relax\detokenize\expandafter{\ExerciseTitle}\relax
		%\marginpar{[Bonus]}
		\else
		\marginpar{\mbox{[\ExerciseTitle]}}
		\fi
		#2
	\end{Exercise}
	\begin{Answer}[ref=#3]
		#4
	\end{Answer}
}
\renewcommand{\exemulticols}[5]{
	\begin{multicols}{2}
	\begin{Exercise}[#1, label=#4]
		\if\relax\detokenize\expandafter{\ExerciseTitle}\relax
		%\marginnote{[Bonus]}
		\else
		\marginnote{\mbox{[\ExerciseTitle]\qquad}}
		\fi
		#2
	\end{Exercise}
	\columnbreak
		#3
	\end{multicols}
	\begin{Answer}[ref=#4]
		#5
	\end{Answer}
}





\SetDate[5/09/2026]

\begin{document}
\fancyhead[C]{\bf Algèbre I --- Série 1}

\subsection*{Automatismes}

\exe{}{
	Soient $f(x) = 3+2x$ et $g(x) = x^2 + 2x - 1$ deux polynômes à coefficients réels.
	Calculer
	\begin{multicols}{4}
	\begin{enumerate}
		\item $f(x)+g(x)$ ;
		\item $f(x) g(x)$ ;
		\item $f(x)^2$ ; et
		\item $f(g(x))$.
	\end{enumerate}
	\end{multicols}
}{exe:sum-product}{
	todo
}

\exe{}{
	Calculer de tête le coefficient
	\begin{multicols}{2}
	\begin{enumerate}
		\item
		en $x$ de $(x+3)(x-4)$ ;
		\item
		en $x^2$ de $(1+x+x^2+x^3)(2x^2 + 3x - 5)$ ; \columnbreak
		\item
		en $x^6$ de $(3-4x-x^3+7x^4 + 2x^5)^2$ ;  et
		\item
		en $x^{50}$ de \mbox{$(1 + x^{25} + x^{50} + x^{100})(1 - x^{25} - x^{50} - x^{100})$}.
	\end{enumerate}
	\end{multicols}
}{exe:coeff-reading}{

}

\subsection*{Exercices}

\exe{}{
	Montrer que le polynôme $x + 1$ est combinaison linéaire des polynômes $x + 3$ et $x+5$.
}{exe:comb-lin-poly}{

}
\exe{}{
	Montrer que le polynôme $x^2 - 5x + 6$ n'est pas combinaison linéaire des polynômes $x - 2$ et $x-3$.
}{exe:not-comb-lin-poly}{

}


\exe{difficulty=1}{
	Soit $S = \{ -1 ; 0 ; 1\}$.
	Considérons l'ensemble $S[x]$, les polynômes à coefficients dans $S$.
	Par définition, le polynôme $x^3 - x$ n'est pas le polynôme nul.
	
	Montrer que la fonction polynomiale $f(x) = x^3 -x$, elle, est la fonction nulle sur $S$.
	
	\emph{Remarque : il est donc important de distinguer polynôme et fonction polynomiale.}
}{exe:fonction-nulle-S}{
	Il s'agit de montrer que $f(x) = 0$ pour tout $x\in S$.
	Ceci se vérifie par le calcul : $f(0) = f(1) = f(-1) = 0$.
}

\exe{}{
	Montrer que la famille de polynômes $\{ x+1 ; x+2 \} \subseteq\R[x]$ est linéairement indépendante sur $\R$.
}{exe:lin-indep1}{
	todo
}

\exe{}{
	Montrer que la famille de polynômes $\{ x^2+4x - 2 ; x+2 ; 50 \} \subseteq\R[x]$ est linéairement indépendante sur $\R$.
}{exe:lin-indep2}{
	todo
}

\exe{}{
	Est-ce que la famille de polynômes $\{ 2x^2+4x - 2 ; x^2 + x + 1 ; x-2 \} \subseteq\R[x]$ est linéairement indépendante sur $\R$ ?
	Démontrer que oui ou fournir une combinaison linéaire nulle et non triviale.
}{exe:lin-indep3}{
	La famille est liée car
		\[ (2x^2 + 4x - 2) - 2(x^2 + x + 1) - 2(x-2) = 0. \]
}




\exe{}{
	Pour chaque propriété, donner un polynôme $f$ à coefficients réels et non nul la vérifiant.
	\begin{multicols}{2}
	\begin{enumerate}
		\item $f$ s'annule en $1$.
		\item $f$ s'annule en $-10$.
		\item $f$ s'annule en $0$.
		\item $f$ s'annule en $-10$ et en $1$.
		\item Les racines de $f$ sont $2, -3, \frac27$, et $0$.
	\end{enumerate}
	\end{multicols}
}{exe:6}{
	Le polynôme identiquement nul $f(x) = 0$ n'est bien sûr pas autorisé sinon l'exercice n'a pas d'intérêt !
	\begin{enumerate}
		\item $f(x) = x-1$ fonctionne ainsi que $f(x) = 3(x-1)$, ou $f(x) = -(x-1) = 1-x$.
		\item $f(x) = x+10$ fonctionne, ainsi que $f(x) = (x+10)(x-1)$, et tous ses multiples.
		\item $f(x) = x$ ou $f(x) = 500x$.
		\item $f(x) = (x+10)(x-1) = x^2 +9x -10$.
		\item $f(x) = (x-2)(x+3)(x-\frac27)x$.
	\end{enumerate}

}



\exe{}{
	Considérons l'équation quadratique $3x^2-7x+4= 0$ pour $x\in\R$.
	Montrer que $1$ est racine de l'équation et en déduire l'autre racine.
}{exe:euclide-poly}{
	todo
}




\exe{}{
	Trouver une racine triviale (évidente) puis factoriser le polynôme $f(x) = 7x^2 - 10x + 3$ pour trouver sa deuxième racine.
}{exe:euclide-poly2}{
	todo
}


\exe{}{
	Considérons le polynôme $f(x) = 21x^3 - 7x^2 - 9x + 3$.
	Montrer, par le calcul, que $\frac13$ est une racine de $f$.
	En déduire les autres racines de $f$.
	
	\emph{Indication : calculer la division euclidienne de $f$ par $3x-1$ plutôt que par $x-\frac13$ pour simplifier les calculs.}
}{exe:euclide-poly3}{
	$f(x) = (3x-1)(7x^2-3)$ donc les racines sont $\frac13$ et $\pm\sqrt{\frac37}$.
}

\exe{}{
	Considérons le polynôme $f(x) = 4x^3 - 6x^2 - 2x + 3$.
	Montrer, par le calcul, que $\frac32$ est une racine de $f$.
	En déduire les autres racines de $f$.
}{exe:8}{
		\[ (2x - 3)(2x^2 - 1) = 4x^3 - 2x -6x^2 + 3. \]
	On utilise que le produit est nul si et seulement si un des deux facteurs est nul.
	Or l'équation $2x^2 - 1 = 0$ est équivalente à $x^2 = \dfrac12$.
	En prenant la racine et en utilisant que $\sqrt{x^2} = |x|$, on obtient
		\[ |x| = \sqrt{\dfrac12} = \dfrac1{\sqrt2}. \]
	Les deux valeurs $\pm \dfrac1{\sqrt2}$ (notation $\pm$ pour \og plus ou moins \fg) sont donc solutions.
	On en déduit que l'ensemble des solutions de \eqref{eq:2} est $\left\{ \dfrac32 ; \pm \dfrac1{\sqrt2} \right\} = \left\{ \dfrac32 ; \dfrac1{\sqrt2} ; - \dfrac1{\sqrt2} \right\}$

}


\exe{difficulty=1}{
	Montrer le théorème de Taylor bis : 
	soit $f\in\R[x]$ un polynôme à coefficients réels de degré $d$ et $a\in\R$ un réel quelconque.
	Alors
		\[ f(x) = \sum_{k=0}^{d} \frac{f^{(k)}(a)}{k!} (x-a)^k.\]
}{exe:taylor-poly-2}{
	todo
}



\exe{}{
	Considérons le polynôme
		\[ f(x) = (x-3)^2 - (x-3)^4 + 3(x-3)^5. \]
	Calculer $f(3), f'(3), f''(3), f'''(3), f^{(\text{iv})}(3)$, et $f^{(\text{v})}(3)$.
}{exe:f-shift-derivatives}{

}


\exe{difficulty=1}{
	Considérons le \emphindex{symbole de Pochhammer} défini par
		\begin{align*}
			(x)_0 = 1, && \et && (x)_n = x(x-1)(x-2)\cdots(x-n+1) \text{ pour $n\in\N^*$, }
		\end{align*}
	ainsi que l'opérateur de différence $\Delta$ qui associe à chaque $f\in\R[x]$ le polynôme $f(x+1)-f(x)$.
	
	Montrer que, pour tout $n\in\N^*$, l'identité suivante est vraie dans $\R[x]$.
		\[ \Delta(x)_n = n\cdot(x)_{n-1} \]
}{exe:Pochhammer}{

}



\exe{difficulty=1}{
	Considérons une famille de polynômes $\{ f_1, f_2, \dots, f_n\} \subseteq \R[x]$ telle que pour tout $i=1, 2, \dots n$, il existe un nombre $\alpha_i\in\R$ qui est racine de $f_j$ pour tout $j\neq i$, mais n'est pas racine de $f_i$.
	Montrer que la famille est linéairement indépendante sur $\R$.
	
	Conclure que la famille $F$ suivante est linéairement indépendante sur $\R$.
		\[ F = \bigset{ (x-1)(x-2)(x-3) ; x(x-2)(x-3) ; x(x-1)(x-3) ; x(x-1)(x-2) } \]
	Est-ce qu'il est possible d'écrire n'importe quel polynôme de degré inférieur ou égal à 3 comme combinaison linéaire des éléments de $F$ ?
}{exe:unique-roots}{
	todo
}



\exe{}{
	Vrai ou faux ?
	\begin{enumerate}%[itemsep=10pt]
		\item
		Si $f|g$ dans $\R[x]$, alors $\deg(f) \leq \deg(g)$.
		\item
		Si $f$ divise $g$ dans $\R[x]$, alors $(3f)$ divise aussi $g$.
		\item
		Soient $f, g\in\R[x]$. Alors $\deg(fg) = \deg(f)+\deg(g)$.
		\item
		Soient $f, g\in\R[x]$. $\deg(f+g) = \max\{\deg(f) ; \deg(g)\}$.
		\item
		Un polynôme constant de $\R[x]$ est de degré nul.
		\item
		$7x$ divise $x$ dans $\R[x]$.
		\item
		$124(x+1)$ divise $2x^5 -x^4 - 3x^3 + x^2 -3x - 4$ dans $\R[x]$.
	\end{enumerate}
}{exe:TF-polynomes}{
	\begin{enumerate}
		\item
		Faux uniquement à cause du polynôme nul $g=0$. 
		Tous les polynômes divisent le polynôme nul. 
		\item
		Vrai.
		\item
		Vrai car le coefficient dominant de $fg$ est le produit des coefficients dominants de $f$ et de $g$.
		\item
		Il faut et il suffit que la somme des coefficients de degré $\max\{\deg(f) ; \deg(g)\}$ soit non nulle.
		Si $g=-f$ par exemple, alors $\deg(f+g) = \minfty$.
		\item
		Faux à cause du polynôme nul de degré $\minfty$ par convention.
		\item
		Vrai car $(7x)\cdot\frac17 = x$.
		\item
		Vrai car $-1$ est racine du deuxième polynôme.
	\end{enumerate}
	
}


\subsection*{Approfondissements}

\exe{difficulty=1}{
	Soit $d\in\N$ pas un carré parfait. 
	Montrer que la famille $\bigset{ 1 ; \sqrt{d} } \subseteq\R$ est linéairement indépendante sur $\Q$.
}{exe:lin-indep4}{
	todo
}

\exe{difficulty=1}{
	Soit $F_1, F_2 \subseteq \R[x]$ deux familles de polynômes.
	Nommons $V_1$ et $V_2$ l'ensemble de toutes les combinaisons linéaires de $F_1$ et $F_2$ respectivement.
	
	Montrer que, pour un polynôme $f$ combinaison linéaire de $F_1 \cup F_2$,
	\begin{align*}
		V_1 \cap V_2 = \{ 0 \} && \iff && f = v_1 + v_2 \text{ se décompose de façon unique avec $v_1\in V_1$ et $v_2\in V_2$}.
	\end{align*}
	Lorsque la décomposition est unique, nous dirons alors que $V_1$ et $V_2$ sont en \emphindex{somme directe}, notée $V_1 \oplus V_2$.
}{exe:direct-sum}{
	todo
}

\exe{}{
	Montrer qu'une fonction polynomiale de $\R[x]$ de degré 3 admet au moins une racine réelle.
}{exe:deg3-root}{
	todo
}

\exe{difficulty=2}{
	Montrer la \emphindex{règle du produit} : pour deux polynômes $f, g\in\R[x]$,
		\[ (fg)' = f'g + g'f. \]
}{exe:deriv-prod}{
	À gauche, le coefficient en $x^n$ de $fg$ est égal à
		\[ \sum_{i+j=n} f_i g_j. \]
	Le coefficient en $x^{n-1}$ de $fg$ est donc égal à
		\[ n\sum_{i+j=n} f_i g_j. \]
	À droite, le coefficient en $x^{n-1}$ de $f'g$ est égal à
		\[ \sum_{i+j=n-1} f'_ig_j = \sum_{i+j=n-1} (i+1)f_{i+1}g_j = \sum_{i+j=n} if_ig_j. \]
	Le coefficient en $x^{n-1}$ de $g'f$ est égal à
		\[ \sum_{i+j=n-1} f_ig'_j = \sum_{i+j=n-1} f_i (j+1)g_{j+1} = \sum_{i+j=n} f_ijg_j. \]
	Le coefficient en $x^{n-1}$ de $f'g + g'f$ est donc égal à
		\[ \sum_{i+j=n} f_ijg_j +  \sum_{i+j=n} if_ig_j = \sum_{i+j=n} (i+j)f_ig_j = n\sum_{i+j=n} f_i g_j. \]
	Les coefficients de $(fg)'$ et $f'g+g'f$ sont donc égaux et les polynômes sont égaux.
}


\exe{difficulty=1}{
	%Soit $f\in C^1(\R)$ une fonction continûment dérivable.
	Soit $f$ une fonction dérivable en 0.
	Montrer que 
		\[ f(x) = f(0) + f'(0)x + g(x), \]
	où $g(x) = o(x)$ dans le sens suivant : 
		\[ \lim_{x \to 0} \frac{g(x)}x = 0. \]
}{exe:taylor-C1}{
	Posons $g(x) = f(x) - f(0) - f'(0)x$.
	Pour $x\neq0$, nous avons
		\[ \frac{g(x)}x = \frac{f(x) - f(0)}{x-0} - f'(0). \]
	Lorsque $x$ tend vers $0$,  $\frac{f(x) - f(0)}{x-0}$ tend vers $f'(0)$ par définition de $f'(0)$.
	Il suit donc que $\frac{g(x)}x$ tend vers 0 comme requis.
	
	Remarque : si $f\in C^2(\R)$, il s'avère que $g(x) = \O(x^2)$, c'est-à-dire que $\frac{g(x)}{x^2}$ reste borné lorsque $x$ tend vers $0$ (ce qui est plus fort que $g=o(x)$).
	Nous aurons en fait 
		\[ f(x) = f(0) + f'(0)x + \frac{f''(0)}2 x^2 + o(x^2). \]
}

\exe{difficulty=1}{
	Le but de cet exercice est de construire explicitement l'ensemble $\R[x]$ dans lequel les monômes $\{ 1 ; x ; x^2 ; \dots \}$ sont linéairement indépendants sur $\R$.
	
	Considérons à cette fin l'ensemble $\hat{\R}^\infty$ des suites réelles ayant seulement un nombre fini de termes non nuls.
	\begin{enumerate}
		\item
		Montrer que toute suite $a\in\hat{\R}^\infty$ prend la forme
			\[ a = (a_0, a_1, \dots, a_d, 0, 0, 0,\dots). \]
	\end{enumerate}
	Nous verrons donc dans la suite un polynôme  $a_0 + a_1 x + \dots + a_d x^d$ comme une suite $a = (a_0, a_1, a_2, \dots, a_d, 0, 0, 0,\dots)$ de $\hat{\R}^\infty$.
	\begin{enumerate}[resume]
		\item
		À quelle suite correspond chaque polynôme suivant ?
		\begin{multicols}{2}
		\begin{enumerate}
			\item Le polynôme nul.
			\item Le polynôme $1$.
			\item Le polynôme $x$.
			\item Le polynôme $x^2$.
		\end{enumerate}
		\end{multicols}
		\item
		Comment définir les multiplication $\alpha a$ et $a \cdot b$ pour $\alpha\in\R$ et $a, b\in\hat{\R}^\infty$ ?
		\item
		Montrer que les suites $\{ e_0 ; e_1 ; \dots ; e_d \}$ sont linéairement indépendantes sur $\R$, où
			\[ e_i = ( 0, 0, \dots, 0, \underbrace{1}_{\text{position $i$}}, 0, 0, \dots ). \]
		Ces suites sont appelées \emphindex{vecteurs canoniques}.
	\end{enumerate}
}{exe:Rx-constr}{

}




%%%%%%%%%%%

\newpage
\fancyhead[C]{\textbf{Solutions}}
\shipoutAnswer

\end{document}
